\section{Initial conditions}
\label{sec:inits}

\subsection{Using restart files}
\label{sec:inits-restarts}

This section describes how to run \NEMOSIcu{} with initial conditions taken from restart files saved by a previous run.
As an example, we shall extend the three-year run of \refsection{sec:inputs-core-ciaf} for another three years.

If you look in the run directory of \refsection{sec:inputs-core-ciaf} (\dirname{EXP02} there), as well as the output files previously listed, you will see several files with names beginning \linebreak\filename{ORCA2_00008760_restart} (assuming the run completed).
These are restart files written after the final time step, which was iteration 8760.

Start by configuring a new run directory in the same way as for the original experiment in \refsection{sec:inputs-core-ciaf}.
You can call this new directory anything you like, but in this example let us suppose it is called \dirname{EXP02-1}.

Now make the following edits in the new run directory.
(As noted previously, when adding new variables to \filename{namelist_cfg} or \filename{namelist_ice_cfg} I generally copy and paste the appropriate lines from \filename{namelist_ref} or \filename{namelist_ice_ref} as a starting point.)

In \verb|namelist_cfg|, make the following changes and additions in namelist group \nmlgrp{namrun}:

\begingroup\footnotesize
\begin{VerbatimNML}[samepage=true]
\nmlpar{nn_it000}        = \nmlval{8761}  \nmlcomment{the old nn_itend + 1}
\nmlpar{nn_itend}        = \nmlval{17520} \nmlcomment{the total number of timesteps in 6 years}
\nmlpar{ln_rstart}       = \nmltrue{}
\nmlpar{nn_rstctl}       = \nmlval{2}
\nmlpar{cn_ocerst_in}    = \nmlstring{ORCA2_00008760_restart}
\nmlpar{cn_ocerst_indir} = \nmlstring{../EXP02} \nmlcomment{(relative) path to the original run directory}
\end{VerbatimNML}
\endgroup

Look at the comments in \filename{namelist_ref} if you want to know what these variables mean.

This is not essential, but doing so will avoid a warning: in namelist group \nmlgrp{namtsd}, set \nmlpar{ln_tsd_init} to \nmlfalse{}.

In \filename{namelist_ice_cfg}, make the following additions in namelist group \nmlgrp{nampar}:

\begin{VerbatimNML}[samepage=true]
\nmlpar{cn_icerst_in}    = \nmlstring{ORCA2_00008760_restart_ice}
\nmlpar{cn_icerst_indir} = \nmlstring{../EXP02}
\end{VerbatimNML}

and in namelist group \nmlgrp{namini} add:

\begin{VerbatimNML}
\nmlpar{nn_iceini_file} = \nmlval{2}
\end{VerbatimNML}

Look at the comments in \filename{namelist_ice_ref} if you want to know what these variables mean.

For this extension experiment, \dirname{EXP02-1}, the job script can be a copy of that for original experiment, \dirname{EXP02}.
If the extension covered a different number of years, the time limit in the job script might need adjustment, but that is not the case in this example.
